%% kapa: exam=1 round=29 year=2018 period=1 form=A subject=real_estate subject_ko=부동산학원론 seq=33 number=113
%% source: Q-Net 감정평가사 기출문제 (cst00309) artlSeq=5202808 file=extracted/1008/5202808/2018년 제29회 감정평가사 1차시험 기출문제/2018년 제29회 감정평가사 1차시험 기출문제/1교시 A형.pdf
%% conversion: kordoc 4.24.51 -> tools/gen1.py (auto); review=proofread-sonnet; flags=none
\begin{question}{113}
우리나라의 부동산가격공시제도에 관한 설명으로 옳은 것은?
\begin{choices}
  \item 다가구주택은 공동주택가격의 공시대상이다.
  \item 개별공시지가의 공시기준일이 6월 1일인 경우도 있다.
  \item 표준주택에 그 주택의 사용ㆍ수익을 제한하는 권리가 설정되어 있을 때에는 이를 반영하여 적정가격을 산정하여야 한다.
  \item 국세 또는 지방세 부과대상이 아닌 단독주택은 개별주택가격을 결정ㆍ공시하지 아니할 수 있다.
  \item 표준지공시지가의 공시권자는 시장ㆍ군수ㆍ구청장이다.
\end{choices}
\end{question}
\answer{4}
